\section{common/numvec.hpp File Reference}
\label{numvec_8hpp}\index{common/numvec.hpp@{common/numvec.hpp}}


\subsection{Detailed Description}


Definition in file {\bf numvec.hpp}.

{\tt \#include \char`\"{}commoninc.hpp\char`\"{}}\par
\subsection*{Classes}
\begin{CompactItemize}
\item 
class {\bf Num\-Vec$<$ F $>$}
\begin{CompactList}\small\item\em Num\-Vec is a vector made up of a size and data. The data can be made up of any type (ints, double, bools, etc.). \item\end{CompactList}\end{CompactItemize}
\subsection*{Typedefs}
\begin{CompactItemize}
\item 
typedef {\bf Num\-Vec}$<$ bool $>$ {\bf Bol\-Num\-Vec}
\item 
typedef {\bf Num\-Vec}$<$ int $>$ {\bf Int\-Num\-Vec}
\item 
typedef {\bf Num\-Vec}$<$ double $>$ {\bf Dbl\-Num\-Vec}
\end{CompactItemize}
\subsection*{Functions}
\begin{CompactItemize}
\item 
template$<$class F$>$ ostream \& {\bf operator$<$$<$} (ostream \&os, const {\bf Num\-Vec}$<$ F $>$ \&vec)
\item 
template$<$class F$>$ void {\bf setvalue} ({\bf Num\-Vec}$<$ F $>$ \&V, F val)
\item 
template$<$class F$>$ void {\bf clear} ({\bf Num\-Vec}$<$ F $>$ \&V)
\end{CompactItemize}


\subsection{Typedef Documentation}
\index{numvec.hpp@{numvec.hpp}!BolNumVec@{BolNumVec}}
\index{BolNumVec@{BolNumVec}!numvec.hpp@{numvec.hpp}}
\subsubsection{\setlength{\rightskip}{0pt plus 5cm}typedef {\bf Num\-Vec}$<$bool$>$ {\bf Bol\-Num\-Vec}}\label{numvec_8hpp_e7b334ad2709d68109d453b8e618c9cd_e7b334ad2709d68109d453b8e618c9cd}


Numvec of bools 

Definition at line 127 of file numvec.hpp.\index{numvec.hpp@{numvec.hpp}!DblNumVec@{DblNumVec}}
\index{DblNumVec@{DblNumVec}!numvec.hpp@{numvec.hpp}}
\subsubsection{\setlength{\rightskip}{0pt plus 5cm}typedef {\bf Num\-Vec}$<$double$>$ {\bf Dbl\-Num\-Vec}}\label{numvec_8hpp_440ba0a3cc36d52c2c5f9b540bed8ac5_440ba0a3cc36d52c2c5f9b540bed8ac5}


Numvec of doubles 

Definition at line 131 of file numvec.hpp.\index{numvec.hpp@{numvec.hpp}!IntNumVec@{IntNumVec}}
\index{IntNumVec@{IntNumVec}!numvec.hpp@{numvec.hpp}}
\subsubsection{\setlength{\rightskip}{0pt plus 5cm}typedef {\bf Num\-Vec}$<$int$>$ {\bf Int\-Num\-Vec}}\label{numvec_8hpp_d02385ca3d87c7e12eab3349354f41ef_d02385ca3d87c7e12eab3349354f41ef}


Numvec of ints 

Definition at line 129 of file numvec.hpp.

\subsection{Function Documentation}
\index{numvec.hpp@{numvec.hpp}!clear@{clear}}
\index{clear@{clear}!numvec.hpp@{numvec.hpp}}
\subsubsection{\setlength{\rightskip}{0pt plus 5cm}template$<$class F$>$ void clear ({\bf Num\-Vec}$<$ F $>$ \& {\em V})\hspace{0.3cm}{\tt  [inline]}}\label{numvec_8hpp_e9a0f533123ffe65bbb091dd24094958_e9a0f533123ffe65bbb091dd24094958}


Clear all values of {\bf Num\-Vec}{\rm (p.\,\pageref{classNumVec})} V 

Definition at line 122 of file numvec.hpp.

References Num\-Vec$<$ F $>$::data(), and Num\-Vec$<$ F $>$::m().

\footnotesize\begin{verbatim}123 {
124   memset(V.data(), 0, V.m()*sizeof(F));
125 }
\end{verbatim}\normalsize 


\index{numvec.hpp@{numvec.hpp}!operator<<@{operator$<$$<$}}
\index{operator<<@{operator$<$$<$}!numvec.hpp@{numvec.hpp}}
\subsubsection{\setlength{\rightskip}{0pt plus 5cm}template$<$class F$>$ ostream\& operator$<$$<$ (ostream \& {\em os}, const {\bf Num\-Vec}$<$ F $>$ \& {\em vec})\hspace{0.3cm}{\tt  [inline]}}\label{numvec_8hpp_36c1ac071bd26869e8467036f9a51dc9_36c1ac071bd26869e8467036f9a51dc9}


output stream operator $<$$<$ 

Definition at line 107 of file numvec.hpp.

\footnotesize\begin{verbatim}108 {
109   os<<vec.m()<<endl;
110   os.setf(ios_base::scientific, ios_base::floatfield);
111   for(int i=0; i<vec.m(); i++)   os<<" "<<vec(i);
112   os<<endl;
113   return os;
114 }
\end{verbatim}\normalsize 


\index{numvec.hpp@{numvec.hpp}!setvalue@{setvalue}}
\index{setvalue@{setvalue}!numvec.hpp@{numvec.hpp}}
\subsubsection{\setlength{\rightskip}{0pt plus 5cm}template$<$class F$>$ void setvalue ({\bf Num\-Vec}$<$ F $>$ \& {\em V}, F {\em val})\hspace{0.3cm}{\tt  [inline]}}\label{numvec_8hpp_bd30a0851e742d36c69575eb38314181_bd30a0851e742d36c69575eb38314181}


Set all members of a {\bf Num\-Vec}{\rm (p.\,\pageref{classNumVec})} V equal to val. That is, V(i) = V.\_\-data[i] = val for all i 

Definition at line 116 of file numvec.hpp.

References Num\-Vec$<$ F $>$::m().

\footnotesize\begin{verbatim}117 {
118   for(int i=0; i<V.m(); i++)
119          V(i) = val;
120 }
\end{verbatim}\normalsize 


